\documentclass[11pt,letterpaper]{article}
\usepackage[T1]{fontenc}
\usepackage{lmodern,microtype}
\usepackage[margin=0.78in,headheight=15pt,headsep=15pt]{geometry}
\usepackage{mathtools,amssymb}
\usepackage{xcolor,booktabs,array,tikz}
\usepackage[most]{tcolorbox}
\usepackage{enumitem,fancyhdr,needspace,hyperref}
\definecolor{ink}{HTML}{162332}
\definecolor{navy}{HTML}{17324D}
\definecolor{teal}{HTML}{167D82}
\definecolor{mist}{HTML}{EDF5F4}
\definecolor{muted}{HTML}{5D6B78}
\color{ink}
\hypersetup{colorlinks=true,linkcolor=teal,urlcolor=teal,pdfauthor={Igor Rivin}}
\setlength{\parindent}{0pt}
\setlength{\parskip}{4pt}
\setlength{\abovedisplayskip}{7pt}
\setlength{\belowdisplayskip}{7pt}
\setlist[enumerate]{label=(\alph*),leftmargin=1.8em,itemsep=3pt,topsep=4pt}
\newcommand{\Z}{\mathbb Z}
\newcommand{\F}{\mathbb F}
\newcommand{\one}{\mathbf{1}}
\newcommand{\id}{\operatorname{id}}
\newcommand{\pts}[1]{\hfill\textcolor{muted}{\small[#1 points]}}
\newcommand{\partpts}[1]{\unskip\nobreak\hspace{.4em}\textcolor{muted}{\small(#1)}}
\newcommand{\edition}{\ifsolutions Instructor solutions\else Student handout\fi}
\newcounter{problem}
\newcommand{\problem}[2]{\par\addvspace{11pt}\Needspace{6\baselineskip}%
  \stepcounter{problem}{\bfseries\color{navy}\theproblem.\ #1}\pts{#2}\par\nobreak}
\newcommand{\topic}[1]{\par\addvspace{9pt}\Needspace{10\baselineskip}%
  {\large\bfseries\color{teal}#1}\par\nobreak}
\newcommand{\studentbreak}{\ifsolutions\else\newpage\fi}
\newcommand{\challenge}[1]{\par\addvspace{12pt}\Needspace{7\baselineskip}%
  {\bfseries\color{teal}Optional challenge: #1}\pts{10 extra}\par\nobreak}
\pagestyle{fancy}
\fancyhf{}
\fancyhead[L]{\small\color{muted}\setname}
\fancyhead[R]{\small\color{muted}Patterns, Proofs, and Primes}
\fancyfoot[L]{\footnotesize\color{muted}\edition}
\fancyfoot[R]{\small\color{teal}\thepage}
\renewcommand{\headrulewidth}{0.3pt}
\NewDocumentEnvironment{solution}{+b}{\ifsolutions
  \begin{tcolorbox}[enhanced,breakable,colback=mist!45,colframe=teal!45,
    boxrule=.5pt,arc=1mm,left=3mm,right=3mm,top=2mm,bottom=2mm,
    before skip=7pt,after skip=7pt,title={Solution},
    title after break={Solution (continued)},colbacktitle=mist,
    coltitle=teal,fonttitle=\small\bfseries]
  \small #1
  \end{tcolorbox}\fi}{}
\newcommand{\handouttitle}[3]{%
  {\LARGE\bfseries\color{navy}\setname}\hfill{\small\color{teal}\edition}\\[4pt]
  {\large #1}\\[4pt]{\small\color{muted}#2}\par
  \medskip\textbf{Preparation.} #3\par
  \ifsolutions\else\medskip
  \textbf{Name:}\ \rule{.46\linewidth}{.4pt}\hfill
  \textbf{Due:}\ \rule{.20\linewidth}{.4pt}\par\fi}

\newcommand{\Q}{\mathbb Q}
\newcommand{\setname}{Midterm review problem set}
\newcommand{\per}[1]{\overline{#1}}
\DeclareMathOperator{\ord}{ord}
\hypersetup{pdftitle={Number Theory: Midterm Review through Finite Fields, Lecture 2}}
\begin{document}
\handouttitle{From Euclid to finite fields}
{16 review problems \textbar\ 160 practice points + 20 optional
\textbar\ Suggested time: 5--7 hours over several sittings}
{Earlier lectures and assignments, ending with finite-fields Lecture 2:
independence, bases, dimension, prime-power cardinalities, and quadratic extensions.}

\textbf{How to use this set.} This is a review bank, not a timed mock exam.
Try each problem before reading a solution; explain every use of a theorem.
Write on separate paper. The point values indicate relative weight.
You may quote results proved in class except when explicitly asked to prove
them. The final two optional problems revisit the earlier enrichment set.
Finite-field cyclicity, matrix counts, Frobenius, and Fibonacci period
arguments from Lectures 3--4 are outside this set's scope.

\textbf{Notation.} $v_p(m)$ is the exponent of the prime $p$ in a nonzero
integer $m$; extend it to nonzero rationals by
$v_p(a/b)=v_p(a)-v_p(b)$. The functions $\varphi,\tau,\sigma$ count units,
count positive divisors, and sum positive divisors, respectively.
The M\"obius function satisfies $\mu(1)=1$, vanishes on integers divisible
by a prime square, and equals $(-1)^r$ on a product of $r$ distinct primes.

\ifsolutions
\textbf{Instructor note.} Each core problem is worth 10 practice points;
subpart allocations are shown. Accept equivalent arguments and award method
credit after carried arithmetic errors. Suggested time is an estimate, not
a classroom-tested completion time. Coverage follows the earlier assignments;
the finite-field cutoff follows the main route of Lecture 2.
\fi

\topic{I. Integers and continued fractions}
\problem{Euclid gives a certificate}{10}
\begin{enumerate}
\item Compute $\gcd(252,165)$ by Euclid and express it as an integer
linear combination of $252$ and $165$. \partpts{5}
\item Find all solutions of $165x\equiv6\pmod{252}$, giving the distinct
residues modulo $252$. Explain why there are exactly that many. \partpts{5}
\end{enumerate}
\begin{solution}
The successive nonzero remainders are $87,78,9,6,3$, so the gcd is $3$.
Back substitution gives $3=19\cdot252-29\cdot165$.
Multiplying by $2$ gives a solution $x=-58$. Dividing the congruence
by $3$ gives $55x\equiv2\pmod{84}$, with a unique solution modulo $84$
because $\gcd(55,84)=1$. Thus $x\equiv26\pmod{84}$, or
$x\equiv26,110,194\pmod{252}$. These three classes exhaust the solutions.
\end{solution}

\problem{What factorization proves}{10}
\begin{enumerate}
\item Prove that $\sqrt{72}$ is irrational by comparing prime exponents
in a hypothetical equality $a^2=72b^2$, with $a,b$ nonzero integers.
\partpts{4}
\item Prove there are infinitely many primes congruent to $3$ modulo $4$.
\emph{Hint:} given a proposed complete list $p_1,\ldots,p_r$, consider
$N=4p_1\cdots p_r-1$. Explain why some prime divisor of $N$ is $3$ modulo $4$;
do not assume $N$ is prime. \partpts{6}
\end{enumerate}
\begin{solution}
In $a^2=2^3 3^2 b^2$, the exponent of $2$ is even on the left and odd
on the right, contradicting unique factorization. For the second part,
the list is nonempty since $3$ is such a prime. The positive odd integer
$N\equiv3\pmod4$ has only odd prime factors. If all were $1\pmod4$,
their product with multiplicities would be $1\pmod4$. Thus some prime
$q\mid N$ is $3\pmod4$. But $N\equiv-1\pmod{p_i}$ for each $i$, so
$q$ is absent from the list, a contradiction.
\end{solution}

\problem{A finite continued fraction remembers B\'ezout}{10}
\begin{enumerate}
\item Find the simple continued fraction of $157/68$ and all of its
convergents. Use the convention that the last digit is at least $2$.
\partpts{4}
\item Multiply the digit matrices
$D(a)=\begin{psmallmatrix}a&1\\1&0\end{psmallmatrix}$.
Use the determinant of their product to find the inverse of $68$ modulo $157$.
\partpts{4}
\item Compute exactly $157/68-30/13$. \partpts{2}
\end{enumerate}
\begin{solution}
Euclid gives $157=2\cdot68+21$, $68=3\cdot21+5$,
$21=4\cdot5+1$, $5=5\cdot1$. Thus $157/68=[2;3,4,5]$,
with convergents $2/1,7/3,30/13,157/68$.
The product is $\begin{psmallmatrix}157&30\\68&13\end{psmallmatrix}$,
of determinant $1$. Consequently $157\cdot13-68\cdot30=1$, so
$68^{-1}\equiv-30\equiv127\pmod{157}$.
Finally $157/68-30/13=(2041-2040)/884=1/884$.
\end{solution}

\studentbreak
\problem{Periodic is not always purely periodic}{10}
For $x_n=(P_n+\sqrt D)/Q_n$ and $a_n=\lfloor x_n\rfloor$, recall
\[
P_{n+1}=a_nQ_n-P_n,\qquad Q_{n+1}=(D-P_{n+1}^2)/Q_n.
\]
\begin{enumerate}
\item Run this recurrence for $\sqrt{11}$ until a state repeats, and
give its continued fraction. Justify the floors with integer bounds.
\partpts{5}
\item Decide which of $\sqrt{11}$, $3+\sqrt{11}$, and
$(3+\sqrt{11})/2$ have purely periodic expansions. Use Galois's
criterion $x>1$ and $-1<x'<0$, where $x'$ is the conjugate.
Read off the purely periodic expansions from (a). \partpts{5}
\end{enumerate}
\begin{solution}
The states $(P_n,Q_n,a_n)$ are
$(0,1,3),(3,2,3),(3,1,6),(3,2,3)$.
Indeed $3<\sqrt{11}<4$, so
$3<(3+\sqrt{11})/2<7/2$ and $6<3+\sqrt{11}<7$.
Thus $\sqrt{11}=[3;\per{3,6}]$.
Its conjugate is $-\sqrt{11}<-1$, so it is not purely periodic.
The other two numbers exceed $1$, and their conjugates lie in $(-1,0)$.
Their expansions are $3+\sqrt{11}=[\per{6,3}]$ and
$(3+\sqrt{11})/2=[\per{3,6}]$.
\end{solution}

\topic{II. Congruences, powers, and arithmetic functions}
\problem{Reassemble residues, and check cancellation}{10}
\begin{enumerate}
\item Solve $x\equiv3\pmod8$, $x\equiv4\pmod9$, and justify uniqueness
modulo $72$. \partpts{5}
\item Find every solution of $x^2\equiv1\pmod{15}$ using reduction
modulo $3$ and $5$. Why does the argument
``$(x-1)(x+1)=0$, so one factor is zero'' fail modulo $15$? \partpts{5}
\end{enumerate}
\begin{solution}
Set $x=3+8k$. Modulo $9$, $3-k\equiv4$, so $k\equiv8$ and
$x\equiv67\pmod{72}$. A difference of two solutions is divisible by
both $8$ and $9$, hence by $72$; adding multiples of $72$ preserves both
congruences. Modulo each prime $3,5$, the roots are $\pm1$.
The four CRT combinations yield $1,4,11,14$ modulo $15$.
Modulo $15$ there are nonzero zero divisors: for $x=4$,
both $x-1=3$ and $x+1=5$ are nonzero, but their product is zero.
\end{solution}

\problem{Fermat and inverse pairing}{10}
\begin{enumerate}
\item Compute $7^{222}\pmod{13}$ efficiently. \partpts{3}
\item Prove Wilson's theorem $(p-1)!\equiv-1\pmod p$ for a prime $p$
by pairing inverses. Account for $p=2$. \partpts{5}
\item The composite number $341$ satisfies $2^{340}\equiv1\pmod{341}$.
Explain why this is consistent with Fermat's little theorem. \partpts{2}
\end{enumerate}
\begin{solution}
Since $222\equiv6\pmod{12}$, Fermat gives $7^{222}\equiv7^6$.
Now $7^2\equiv10$, $7^3\equiv5$, and $7^6\equiv25\equiv12$.
For odd $p$, each nonzero residue has a unique inverse. The self-inverse
ones solve $(a-1)(a+1)=0$ in $\F_p$, so are exactly $1,-1$.
All other factors pair to $1$, leaving product $-1$.
For $p=2$, $1!=1\equiv-1$. Fermat gives a necessary condition for
primality, not a sufficient one; its converse is false.
\end{solution}

\problem{Orders explain exponent reduction}{10}
\begin{enumerate}
\item Find the orders of $2$ modulo $5$, modulo $7$, and modulo $35$.
Justify the relation between the three orders. \partpts{4}
\item Find $\varphi(35)$ and compute $2^{1000}\pmod{35}$. \partpts{3}
\item Explain how Lagrange's theorem applied to $(\Z/n\Z)^\times$
proves Euler's theorem. State the hypothesis on the base. \partpts{3}
\end{enumerate}
\begin{solution}
The powers modulo $5$ give order $4$; modulo $7$, $2,4,1$ give order $3$.
The congruence $2^k\equiv1\pmod{35}$ holds precisely when it holds
modulo both $5$ and $7$, or when $4\mid k$ and $3\mid k$.
Thus the order is $12$. Also $\varphi(35)=4\cdot6=24$ and
$1000\equiv4\pmod{12}$, so $2^{1000}\equiv16\pmod{35}$.
For $\gcd(a,n)=1$, the residue of $a$ is in the unit group of size
$\varphi(n)$. Its cyclic subgroup has size $\ord_n(a)$, which divides
$\varphi(n)$ by Lagrange. Hence $a^{\varphi(n)}\equiv1\pmod n$.
\end{solution}

\problem{Work one prime at a time}{10}
\begin{enumerate}
\item Factor $360$ and compute $\tau(360),\sigma(360),\varphi(360)$,
$\mu(360)$, and $\mu(30)$. Show the prime-power calculations. \partpts{5}
\item Prove that if $f$ is multiplicative, then
$F(n)=\sum_{d\mid n}f(d)$ is multiplicative. Explain precisely where
coprimality is used. \partpts{5}
\end{enumerate}
\begin{solution}
Since $360=2^3 3^2 5$,
\[
\tau(360)=4\cdot3\cdot2=24,\quad
\sigma(360)=15\cdot13\cdot6=1170,\quad
\varphi(360)=360(1-1/2)(1-1/3)(1-1/5)=96.
\]
Also $\mu(360)=0$ and $\mu(30)=-1$.
For $\gcd(m,n)=1$, unique factorization gives a bijection
between $d\mid mn$ and pairs $a\mid m,b\mid n$, via $d=ab$.
Every such pair is coprime, so $f(ab)=f(a)f(b)$. Therefore
$F(mn)=\sum_{a\mid m}\sum_{b\mid n}f(a)f(b)=F(m)F(n)$.
Also $F(1)=f(1)=1$, as required by the normalized definition.
\end{solution}

\studentbreak
\problem{Invert divisor sums and count exact periods}{10}
\begin{enumerate}
\item Suppose $\sum_{d\mid n}f(d)=n^2$. Use M\"obius inversion to
give a formula for $f(n)$, and evaluate $f(12)$. \partpts{4}
\item Let $A(n)$ count labeled binary words of length $n$ with least
rotational period exactly $n$. Explain why
$2^n=\sum_{d\mid n}A(d)$. Find $A(6)$ and the number of binary necklaces
of length $6$ with least period $6$, where rotations are identified.
Explain the final division. \partpts{6}
\end{enumerate}
\begin{solution}
Inversion gives $f(n)=\sum_{d\mid n}\mu(d)(n/d)^2$.
At $12$, only $d=1,2,3,6$ contribute, giving $144-36-16+4=96$.
Every length-$n$ word has a unique least period $d\mid n$ and is obtained
by repeating a unique length-$d$ word of least period $d$. Thus the
divisor sum holds. Inverting it gives
$A(6)=2^6-2^3-2^2+2=54$.
A word of least period $6$ has six distinct rotations, so there are
$54/6=9$ such necklaces. Words of smaller period would have smaller
rotation classes, which is why the least-period restriction matters.
\end{solution}

\topic{III. Valuations, primes, and $p$-adic distance}
\problem{Count prime factors without expanding factorials}{10}
\begin{enumerate}
\item State Legendre's formula. Use it to find $v_2(40!)$ and the
number of trailing decimal zeros of $40!$. \partpts{3}
\item Put $B_n=\binom{2n}{n}$. Prove that each summand in
\[
v_p(B_n)=\sum_{j\ge1}\left(
\left\lfloor\frac{2n}{p^j}\right\rfloor-
2\left\lfloor\frac{n}{p^j}\right\rfloor\right)
\]
is $0$ or $1$, and deduce $B_n\mid\operatorname{lcm}(1,\ldots,2n)$.
\partpts{4}
\item For prime $p$, prove $p\mid\binom pk$ when $0<k<p$.
Deduce $(X+Y)^p=X^p+Y^p$ as a polynomial identity in characteristic $p$.
\partpts{3}
\end{enumerate}
\begin{solution}
Legendre gives $v_p(n!)=\sum_{j\ge1}\lfloor n/p^j\rfloor$.
Thus $v_2(40!)=20+10+5+2+1=38$, while $v_5(40!)=8+1=9$;
there are $\min(38,9)=9$ trailing zeros.
Writing $n/p^j=q+t$, $0\le t<1$, makes each summand $\lfloor2t\rfloor$,
which is $0$ or $1$. Only powers $p^j\le2n$ can contribute, so
$v_p(B_n)\le\max\{j:p^j\le2n\}$, the valuation of the lcm (zero
if $p>2n$). Compare valuations for all primes.
Finally $p!$ has one factor $p$ and neither $k!$ nor $(p-k)!$ has any.
Thus $v_p\binom pk=1$. The binomial theorem in characteristic $p$
loses all intermediate coefficients, proving the polynomial identity.
\end{solution}

\problem{Recover the scale of the primes}{10}
Let $\pi(x)=\#\{p\le x:p\text{ prime}\}$ and
$\vartheta(x)=\sum_{p\le x}\log p$, with natural logarithms.
\begin{enumerate}
\item Explain why the central binomial coefficient is the largest in
its row, and prove $4^n/(2n+1)\le B_n<4^n$ for $n\ge1$. \partpts{3}
\item Use Problem 10 to prove \partpts{4}
\[
\pi(2n)\ge\frac{2n\log2-\log(2n+1)}{\log(2n)}.
\]
\item You may use the lecture estimate $\vartheta(x)\le Ax$ for a fixed
constant $A>0$ and all $x\ge2$. Show
$\pi(x)\le\sqrt x+2Ax/\log x$.
Explain why these estimates give $\pi(x)\asymp x/\log x$ but do not
give $\pi(x)\sim x/\log x$. \partpts{3}
\end{enumerate}
\begin{solution}
Consecutive coefficients have ratio $(2n-k)/(k+1)$, so they increase
to the center and then decrease. Their sum is $4^n$; the largest is at
least the average and strictly less than the sum.
By Problem 10, $B_n$ divides $L=\operatorname{lcm}(1,\ldots,2n)$.
For each prime $p\le2n$, its factor in $L$ is at most $2n$, whence
$\log B_n\le\log L\le\pi(2n)\log(2n)$.
Combine this with the lower central bound.
There are at most $\sqrt x$ primes at most $\sqrt x$. Each larger prime
up to $x$ contributes at least $\tfrac12\log x$ to $\vartheta(x)$,
giving the stated upper bound. For general sufficiently large $x$,
set $n=\lfloor x/2\rfloor$ in the lower bound and use
$\pi(x)\ge\pi(2n)$. The resulting lower ratio tends to $\log2>0$;
also $\sqrt x=o(x/\log x)$. Thus fixed positive lower and upper
constants work eventually. Such bounds do not identify the limiting
ratio as $1$, or even prove that a limiting ratio exists.
\end{solution}

\problem{Close in a different arithmetic}{10}
For nonzero rational $r$, define $|r|_3=3^{-v_3(r)}$; set $|0|_3=0$
and $d_3(x,y)=|x-y|_3$.
\begin{enumerate}
\item Compute $|27/20|_3$ and $d_3(1/4,7/4)$. \partpts{2}
\item Prove $|x+y|_3\le\max(|x|_3,|y|_3)$ for rational $x,y$.
Prove equality if $|x|_3\ne|y|_3$. \partpts{4}
\item Find the $3$-adic limit in $\Q$ of $S_n=\sum_{j=0}^n3^j$ by
computing the distance to your proposed limit. Does $S_n$ have a finite
real limit? \partpts{4}
\end{enumerate}
\begin{solution}
The values are $1/27$ and $|-3/2|_3=1/3$.
For nonzero $x,y$, write $x=3^a u$, $y=3^b v$, where $u,v$ are
rationals with numerator and denominator prime to $3$, and suppose $a\le b$.
Then $x+y=3^a(u+3^{b-a}v)$ has valuation at least $a$ (or is zero).
If $a<b$, putting the parenthesis over a denominator prime to $3$
shows its numerator is not divisible by $3$, so the valuation is exactly
$a$. Cases involving zero are immediate. Taking $3$ to the negative
valuation proves both assertions.
The geometric sum is $S_n=(3^{n+1}-1)/2$, so the proposed limit is
$-1/2$ and $d_3(S_n,-1/2)=|3^{n+1}/2|_3=3^{-(n+1)}\to0$.
As real numbers these positive sums grow without bound.
\end{solution}

\studentbreak
\topic{IV. Finite fields: through Lecture 2}
\problem{Units, zero divisors, and characteristic}{10}
\begin{enumerate}
\item List the units of $\Z/12\Z$ and exhibit two nonzero residues
whose product is zero. \partpts{2}
\item Prove that a finite field has prime characteristic. Explain
how its multiples of $1$ form a copy of $\F_p$. \partpts{4}
\item A finite commutative ring $R$ has $1\ne0$ and no nonzero zero
divisors. Prove it is a field by considering multiplication by a fixed
nonzero element. Where is finiteness used? \partpts{4}
\end{enumerate}
\begin{solution}
The units are $1,5,7,11$; for instance $2\cdot6=0$ modulo $12$.
In a finite field, two distinct nonnegative multiples of $1$ coincide,
so a positive multiple is zero. Let $m$ be the least positive such
integer. If $m=ab$ with $1<a,b<m$, then
$(a\cdot1)(b\cdot1)=0$ has two nonzero factors, impossible in a field.
Thus $m=p$ is prime. The map $\Z\to K$, $a\mapsto a\cdot1$, has
kernel $p\Z$ and identifies its image with $\F_p$.
For fixed $a\ne0$ in $R$, $ax=ay$ implies $a(x-y)=0$, hence $x=y$.
Multiplication by $a$ is injective, and finiteness makes it surjective.
In particular, $ab=1$ for some $b\in R$. This holds for every nonzero $a$.
\end{solution}

\problem{Why dimension is well defined}{10}
\begin{enumerate}
\item In $\F_5^2$, show that $u=(1,2)$ and $v=(2,1)$ form a basis.
Find the coordinates of $(3,4)$ in this basis. \partpts{3}
\item Let $v_1,\ldots,v_r$ be independent and let $w_1,\ldots,w_m$
span a vector space over a field. Reproduce the exchange argument
proving $r\le m$: after inserting the first $k$ independent vectors,
why can the next insertion remove an old $w$ rather than a previously
inserted $v$? Deduce that any two finite bases have the same length.
\partpts{7}
\end{enumerate}
\begin{solution}
If $au+bv=0$, then $a+2b=0$ and $2a+b=0$, hence $-3b=0$ in $\F_5$,
so $b=a=0$. The map $(a,b)\mapsto au+bv$ is injective on a set of
$25$ elements into $\F_5^2$, also of size $25$; it is onto, so this is
a basis. Also $(3,4)=4v$, giving coordinates $(0,4)$.
After $k$ exchanges, the list $v_1,\ldots,v_k$ and $m-k$ old vectors
still spans. Write $v_{k+1}$ as a combination of this list. Some old
vector has a nonzero coefficient, since otherwise $v_{k+1}$ would lie
in the span of $v_1,\ldots,v_k$, contrary to independence. Solve for
that old vector by dividing by its coefficient; replacing it by
$v_{k+1}$ preserves the span. If $r>m$, after $m$ insertions the
spanning list is $v_1,\ldots,v_m$, forcing $v_{m+1}$ into its span,
a contradiction. Thus $r\le m$. Apply the inequality in both directions
to two bases to obtain equality of their lengths.
\end{solution}

\problem{Count coordinates, not residue classes}{10}
\begin{enumerate}
\item Prove every finite field $K$ has $p^d$ elements for some prime
$p$ and integer $d\ge1$. Explain why a basis over $\F_p$ exists and
why counting its coordinates counts each element exactly once. \partpts{6}
\item If a field has $81$ elements, determine its characteristic and its
dimension over its prime field. Can a field have $12$ elements?
Does the prime-power argument by itself construct a field of every
prime-power size? \partpts{4}
\end{enumerate}
\begin{solution}
By Problem 13, $K$ contains its prime field $\F_p$ and is a vector
space over it using field addition and multiplication by scalars in
$\F_p$. Start with the empty independent list and repeatedly add a
vector outside its span. Such an addition preserves independence:
a relation with nonzero coefficient of the new vector would put it
in the old span. Each addition strictly enlarges the span, so finiteness
forces the procedure to stop at a spanning independent list, a basis.
If it has length $d$, spanning gives existence and independence gives
uniqueness of coordinates in $\F_p^d$. Thus $|K|=p^d$; since $1\ne0$,
$d\ge1$. For $81=3^4$, the characteristic is $3$ and dimension $4$.
The integer $12$ is not a prime power, so no field has that cardinality.
The argument proves necessity; a general existence theorem has not
been proved by this counting argument.
\end{solution}

\problem{Build a field with twenty-five elements}{10}
Use formal pairs $a+bs$ with $a,b\in\F_5$, and the relation $s^2=2$.
Addition is coefficientwise and multiplication is distributive with
$s^2$ replaced by $2$. You may use the resulting commutative ring laws.
\begin{enumerate}
\item List the squares in $\F_5$. Prove that every nonzero $a+bs$
has inverse $(a-bs)/(a^2-2b^2)$, explaining why the denominator is
nonzero. \partpts{4}
\item Compute $(1+s)(2+3s)$ and $(1+s)^{-1}$, and check the inverse
by multiplication. \partpts{2}
\item Give a basis over $\F_5$, the cardinality, and the characteristic
of this field. Explain why it differs from $\Z/25\Z$. \partpts{2}
\item If we instead imposed $t^2=4$ on formal pairs over $\F_5$,
would the resulting ring be a field? Give explicit evidence. \partpts{2}
\end{enumerate}
\begin{solution}
The squares are $0,1,4$, so $2$ is nonsquare. The product of $a+bs$
with $a-bs$ is $a^2-2b^2$. If $b=0$ and $a\ne0$, this is nonzero.
If $b\ne0$, its vanishing would give $(a/b)^2=2$, a contradiction.
Division by this nonzero scalar gives the inverse.
The product is $2+5s+3s^2=8=3$ in $\F_5$.
The inverse is $(1-s)/(1-2)=4+s$; indeed
$(1+s)(4+s)=6+5s=1$.
The basis is $1,s$, since formal pairs have unique coefficients.
The field has $25$ elements and characteristic $5$.
The ring $\Z/25\Z$ has characteristic $25$ and nonzero zero divisor $5$.
For the alternative construction, $(t-2)(t+2)=t^2-4=0$, but both
factors are nonzero formal pairs, so this ring is not a field.
\end{solution}

\studentbreak
\topic{Optional review of earlier enrichment}
These two problems review the assignment explicitly labeled optional
enrichment; they do not introduce a later finite-fields lecture.

\challenge{Inclusion--exclusion as inversion}
In a group of $40$ people, the numbers in activities $A,B,C$ are
$20,18,15$; the pairwise intersection counts are $9,7,6$, and the triple
intersection count is $3$. Pairwise counts include the triple intersection.
Find the number in exactly one activity, exactly two activities, and no
activities. Explain the alternating signs as recovering exact counts
from counts that include additional memberships.
\begin{solution}
The counts in only $A,B,C$ are respectively
$20-9-7+3=7$, $18-9-6+3=6$, and $15-7-6+3=5$.
Thus exactly one has count $18$. Exactly two has count
$(9-3)+(7-3)+(6-3)=13$. The union has
$20+18+15-9-7-6+3=34$, so none has count $6$.
For instance, subtracting $A\cap B$ and $A\cap C$ from $A$ subtracts
the triple intersection twice, so it must be added back once.
Suggested credit: exact single counts 4; exact double counts 2; none 2;
explanation of inversion 2.
\end{solution}

\challenge{Interpolation and a finite Fourier transform}
In $\F_5$, let $\omega=2$, whose powers are $1,2,4,3$.
A polynomial $f$ of degree at most $3$ satisfies
$f(1)=1$, $f(2)=0$, $f(4)=0$, $f(3)=0$.
Recover $f$ using either Lagrange interpolation or
\[
c_k=4^{-1}\sum_{j=0}^3f(\omega^j)\omega^{-jk},
\qquad f(X)=\sum_{k=0}^3c_kX^k.
\]
Check all four values. Why is the degree restriction necessary for
uniqueness? Give another polynomial of higher degree with the same values.
\begin{solution}
Only $j=0$ contributes, so each $c_k=4^{-1}=4$ in $\F_5$.
Thus $f(X)=4(1+X+X^2+X^3)$.
At $1$ this gives $16=1$; at $2,4,3$, respectively, the parenthesized
sums are $1+2+4+3$, $1+4+1+4$, and $1+3+4+2$, all zero modulo $5$.
Two degree-at-most-$3$ polynomials agreeing at four distinct points
have difference zero by the polynomial root bound. Without that
restriction, $f(X)+(X^4-1)$ has the same four values and is different.
Suggested credit: coefficients 4; checks 2; uniqueness argument 2;
counterexample 2.
\end{solution}
\end{document}
